%
% File: chap06.tex
% Author: Solal Rapaport
% Description: Conclusion and Perspectives. Social-engineering thread is dropped from this thesis.
%
\chapter{Conclusion and Perspectives}
\label{chap:conclusion}

This concluding chapter synthesises what the three empirical chapters established about reference stability, commit-graph mutability, and recoverable removal of sensitive files (\Cref{chap:tags,chap:histories,chap:secrets}).
Git protects the identity of stored objects, but it does not guarantee that a familiar tag, branch, or advertised history will continue to designate those objects (\Cref{sec:intro-problem}).
The sections below focus on what becomes visible only when those measurements are considered together.

% \section{Summary of the Thesis}
% \label{sec:concl-summary}

% \subsection{Mutable Tags and Release Integrity}
% \label{subsec:concl-mutable-tags}

% \Cref{chap:tags} compared successive archived states of \OriginsWithTagsShort tag-bearing repositories over nine years~\cite{rapaport2026tagalterations}.
% It found observable tag alterations in \AlteredRepositoriesPercent of that corpus and showed that \PercentMovesContentChange of moves changed the source tree behind an unchanged version name.
% The Nixpkgs and \texttt{tj-actions/\allowbreak changed-files} cases demonstrate why this distinction matters for reproducibility and supply-chain execution, while the chapter's discussion develops the tag-specific implications.

% \subsection{Altered Histories and Repository Trust}
% \label{subsec:concl-altered-histories}

% \Cref{chap:histories} extended the comparison from tag targets to reachable commit graphs and detected \AlteredCommitsShort observable events across \PercentOriginsAltered of the examined origins~\cite{DBLP:conf/kbse/RapaportPTZ25}.
% Its root-cause analysis separates propagated identifier changes from direct modifications and distinguishes content, metadata, and branch-structure alterations.
% \emph{GitHistorian} makes this archival comparison queryable for individual repositories~\cite{replication-package-altered-histories}.

% % Social-engineering subsection intentionally omitted (scope decision in THESIS_PLAN.md).

% \subsection{Secret Removal and the Limits of Repository-Level Remediation}
% \label{subsec:concl-secret-removal}

% \Cref{chap:secrets} recovered and scanned filename-selected files removed through history rewriting.
% Among \num{281728} downloaded, deduplicated blobs, \SI{9.18}{\percent} produced at least one detector finding; \SI{86.4}{\percent} of private-key blobs flagged by gitleaks then passed offline syntactic validation.
% These are detection and Tier~A results, not an active-credential rate: provider-assisted or explicitly authorised checks remain necessary to establish whether any credential still grants access.

\section{Cross-Study Lessons for Repository Integrity}
\label{sec:concl-synthesis}

The preceding studies targeted reference stability, commit-graph mutability, and recoverable removal of sensitive files.
The three questions posed in \Cref{sec:intro-rqs} receive direct empirical answers in \Cref{chap:tags,chap:histories,chap:secrets}; this section distils the implications that emerge only when those measurements are viewed as parts of one integrity problem.

\subsection{Observable Mutability Is Not Self-Explaining}
\label{subsec:concl-mutability-norm}

The tag and history studies detect different events, but they support one joint conclusion: observable alteration is not confined to repositories with little visibility.
Within both studied corpora, the probability of observing at least one alteration rises with GitHub star count, even though low-visibility repositories dominate absolute volume.
Visibility therefore does not supply a stability guarantee.

That correlation does not identify a cause.
Popular projects may receive more releases, maintenance, automation, scrutiny, and archival visits; the observational data cannot separate those mechanisms or infer whether any individual alteration is corrective, accidental, or adversarial.
The supported conclusion is narrower: alterations are observable at ecosystem scale, including in visible projects, and their meaning requires classification and context.

\subsection{Independent Archival as a Prerequisite for Integrity Measurement}
\label{subsec:concl-archive-instrument}

All three studies depend on evidence that outlives the state currently advertised by a forge.
A tag move is visible only if both targets were retained; a history rewrite requires both commit graphs; and a removed file can be examined only if its former content remains reachable.
\SWH supplies that evidence for the archive-scale measurements above.
Its role here goes beyond preservation: it acts as an independent integrity witness whose repeated observations make overwritten states comparable.

Dependence on episodic archival observation also changes how negative results must be read.
An archive can demonstrate an alteration that it observed, but no observation between two visits can prove that no short-lived alteration occurred.
Integrity measurement thus depends jointly on persistent identifiers, archival coverage, visit frequency, and scalable access to the resulting graph.

\subsection{The Remediation Gap: Repository Operation vs.\ External Consequence}
\label{subsec:concl-remediation-gap}

Taken together, the studies distinguish a repository operation from the consequence it is intended to address.
Moving a tag changes what future consumers resolve, but it does not change the source hash expected by an existing package definition.
Rewriting a history changes the advertised graph, but it does not alter what an auditor, archive, fork, or earlier consumer already observed.
Removing a credential from that graph changes repository visibility, but only provider-side revocation changes the authority granted by the credential.

These are instances of one general rule: remediation must operate at the layer where the consequence persists.
Repository mutation can repair the current presentation of a project; it cannot by itself revise an external build record, retract copied evidence, settle the status of an earlier licence, or revoke external authority~\cite{torresarias2016gitmetadata}.
The distinction between repository operations and the external consequences they leave behind is the principal security lesson obtained by connecting the three studies.

\subsection{A Common Architecture for Automatic Detection}
\label{subsec:concl-contrib-synthesis}

The three studies also instantiate one detection architecture across progressively deeper layers of repository state.
At the tag and history layers, comparing successive snapshots is enough to detect and classify changes.
The secret-removal study cannot stop there: the archive records too many file removals for exhaustive inspection, so the method must retrieve and scan a filtered subset before validation can support any claim about sensitive material.
Cheap operations must precede expensive ones, making automatic analysis manageable across hundreds of millions of origins and billions of revisions.

The architecture deliberately separates detection from interpretation.
It can establish that an archived state changed and identify which layer changed, but it does not label an attack or infer intent.
This is a reusable methodological result: at this scale, analysis must run in sequence, with cheaper comparisons screening candidates before costlier retrieval, scanning, or validation, and each step limited to the claims its evidence actually establishes.

\section{A Thesis-Level Assurance Model}
\label{sec:concl-implications}

The cross-study lessons are descriptive: they report what the measurements show, but not what actors must do in response.
This thesis' \emph{assurance model} names the minimal conditions under which repository integrity could be demonstrated rather than inferred from a forge's current view.
This section proposes three requirements for platforms, maintainers, and downstream consumers: preserve mutation evidence, bind consumption to verifiable identity, and close incidents at the correct authority.

\subsection{Preserve Evidence at the Mutation Boundary}
\label{subsec:concl-impl-platforms}

The first requirement is to preserve evidence when references or histories change.
Because forges process every reference update, they can record reference creation, retargeting, deletion, and force-pushes before prior targets disappear from their interfaces.
Append-only, hash-based event logs could preserve enough evidence for later comparison without republishing content deliberately removed from the visible history.
Independent archives remain necessary as external witnesses, while platform logs could reduce detection latency and make integrity evidence available to affected consumers.

\subsection{Bind Consumption Claims to Verifiable Identity}
\label{subsec:concl-impl-developers}

The second requirement is to bind consumption claims to verifiable content identity.
A human-readable reference is useful for discovery and communication, but a build, provenance statement, or audit record that requires reproducibility should also retain and verify the resolved content identifier.
Platforms and maintainers can strengthen this binding through declared immutable-release policies, signed statements, and transparent notifications when a published reference changes; consumers should not infer those guarantees from naming conventions alone.
This principle supports both security, by exposing substitution, and operational safety, by making the exact input reproducible.

\subsection{Close Incidents at the Layer of Authority}
\label{subsec:concl-impl-downstream}

The third requirement is to remediate each incident where its effects still live, not only in the repository view.
A repository edit is enough when the only remaining problem is what the project currently displays.
When consequences have spread elsewhere, other steps are required: update or reject affected build records, keep historical evidence for audit questions, and revoke or rotate leaked credentials at the issuing provider.
Treating a clean current branch as proof that the underlying incident is closed would mistake repository cleanup for full remediation.

\section{Limitations of the Thesis}
\label{sec:concl-limitations}

The chapter-specific validity analyses detail local threats (\Cref{sec:tags-validity,sec:hist-validity,sec:sec-validity}); three limits apply across the thesis.

\textbf{The measurements are lower bounds on observable events.}
Snapshot comparison misses alterations made and reverted between visits, merges multiple intervening events, and excludes repositories without repeated archival coverage.
The archive can prove observed differences, not continuous stability.

\textbf{The evidence is structural, not causal.}
Repository states do not reveal maintainer intent, and the same operation may reflect maintenance, automation, accident, or attack.
The analyses classify what changed; case studies and external validation are required before stronger security interpretations.
The corpora are also restricted to Git, whose object and reference semantics shape both the detection methods and the resulting claims.

\textbf{Credential activity remains unresolved.}
\Cref{chap:secrets} establishes recoverability, scanner detection, and offline syntactic validity for selected archived files, but it does not establish whether detected credentials remain active or whether maintainers rotated them.
That question remains open until provider-assisted or explicitly authorised validation is completed.

\section{Perspectives}
\label{sec:concl-perspectives}

The cross-study lessons and assurance model above point to capabilities that do not yet exist at ecosystem scale.
The following subsections outline three directions for turning integrity evidence into durable practice.

\subsection{Longitudinal Integrity Infrastructure}
\label{subsec:concl-future-observability}

The retrospective analyses could become longitudinal integrity infrastructure.
One path is to update alteration indices as new archival visits arrive; another is to combine those durable observations with forge-native reference and force-push event streams for lower-latency alerts.
Publishing precomputed indices alongside graph releases would let registries, auditors, and research tools query mutability without reproducing the full computational pipeline.

A broader comparative programme should also test which conclusions are specific to Git.
Applying equivalent before-and-after observations to other version control systems would require system-specific definitions, but it would distinguish consequences of Git's reference semantics from more general maintenance practices.

\subsection{Explicit Integrity Contracts}
\label{subsec:concl-future-detection}

Current tooling often leaves the stability expected from a version name implicit.
Future work should make that expectation machine-readable: a producer could declare whether a reference is mutable, append-only, or permanently bound; a consumer could declare whether it requires exact content identity, provenance continuity, or only the latest compatible source.
Package managers, continuous integration (CI) systems, software bills of materials, and SLSA-style provenance could then verify the appropriate contract instead of assigning one meaning to every reference.

Such contracts also create an evaluation framework for platform mechanisms.
Signed references, immutable releases, transparency logs, and content-addressed lockfiles should be compared by the evidence they preserve and by the failures they prevent, rather than treated as interchangeable integrity features.

\subsection{From Detection to Accountable Remediation}
\label{subsec:concl-future-developer-behavior}

Finding an alteration is only the first step.
Future alert systems should report what changed, the identifiers involved before and after, when the observation was made, and the limits of that evidence, so reviewers can respond proportionately rather than treat every event as an attack.
Ranking by dependency impact or branch role could further focus attention without replacing judgment.

A hosted successor to \emph{GitHistorian} could expose this information to CI, registries, software bills of materials, audit tools, and researchers through one interface.

The harder open problem is remediation.
When a system surfaces a moved tag, a rewritten history, or a removed credential, do organisations update affected builds, and revoke credentials at the provider, or only clean the current repository view?
Qualitative studies of maintainer and consumer practice, combined with longitudinal event data, could measure that gap between alert and genuine fix~\cite{DBLP:conf/secdev/BasakNRW22}.

\section{Final Remarks}
\label{sec:concl-final}

The thesis shows that the properties needed for supply-chain security and operational safety cannot be inferred from the current forge state alone.
However, it does not aim at showing that every repository alteration is harmful or that current forge state is useless
Git protects object identity, while consumers frequently depend on continuity of names, histories, and external authority that Git does not enforce.
Those claims require complementary evidence: verified content identifiers at consumption time, records that preserve superseded states, and remediation performed by the authority that controls the remaining consequence.

In a software supply chain built on mutable references, integrity must be demonstrated by evidence that outlasts the last push.
